\documentclass[10pt,a4paper]{amsart}
\usepackage[margin=19mm]{geometry}
\usepackage{amsmath,amssymb,mathtools}
\usepackage[scr=rsfs,cal=euler]{mathalfa}
\usepackage{xfrac}
\usepackage[unicode,hidelinks,bookmarks=false]{hyperref}
\newcommand{\ie}{i.e.\ }
\newcommand{\cf}{cf.\ }
\newcommand{\resp}{respectively\ }
\newcommand{\Subsection}{\subsection}
\providecommand{\nobreakdash}{\nobreak\hbox{-}}
\setlength{\emergencystretch}{2em}
\multlinegap0pt
\usepackage{bm}
\usepackage{bbm}
\usepackage{dsfont}
\usepackage{xspace}
\usepackage{cancel}
\usepackage{mathtools}
\usepackage[shortlabels]{enumitem}
\usepackage{amsthm}
\makeatletter
\@ifundefined{theorem}{\newtheorem{theorem}{Theorem}}{}
\@ifundefined{cor}{\newtheorem{cor}[theorem]{Corollary}}{}
\@ifundefined{lemma}{\newtheorem{lemma}[theorem]{Lemma}}{}
\makeatother
\newcommand{\EndsInf}[1]{\mathcal{E}_{\infty}\left(#1\right)}
\newcommand{\HH}{\mathbb{H}}
\newcommand{\Hypers}{\mathcal{H}}
\newcommand{\MM}{\mathbb{M}}
\newcommand{\NN}{\mathbb{N}}
\title[Topology and dynamics of unimodular random hyperbolic manifo]{Topology and dynamics of unimodular random hyperbolic manifolds}
\author{Ilya Gekhtman, Nir Lazarovich, Arie Levit, Asaf Nachmias}
\date{Source: arXiv:2607.25065v1 (CC BY 4.0)}
\begin{document}
\maketitle

\noindent\emph{Source and licence.} Topology and dynamics of unimodular random hyperbolic manifolds. Version arXiv:2607.25065v1 (CC BY 4.0), distributed under the Creative Commons Attribution 4.0 International licence, as recorded in the arXiv licence metadata for this exact version and shown on the abstract page https://arxiv.org/abs/2607.25065v1. Full rights notice: licenses/arxiv-2607.25065-CC-BY-4.0.txt.\par\smallskip

\noindent\textit{Editorial scope.} Counted unit: the Theorem whose source label is \texttt{thm:random walk does not converge to a single end} in the article (source0.tex lines 1730--1733, source label \texttt{thm:random walk does not converge to a single end}), delivered together with the article's own complete original proof of it (source0.tex lines 1737--1790, one \texttt{proof} environment of 54 lines). No mathematics is written, repaired or re-derived: statement and proof are the article's own text line for line. The proof is detached from the statement in the source: the article states the result at lines 1730--1733 and prints its proof at lines 1737--1790, and the proof environment's own header names the statement, so the association is the article's own. Required context retained uncounted: lemma:MTP (lines 425--431); cor:thick part is infinitely ended (lines 1616--1619). Every definition, display and label of those lines is retained unchanged. Omitted from this package and not redistributed: 1--424, 432--1615, 1620--1729, 1734--1736, 1791--1959. Nothing inside a retained excerpt is omitted or altered: every retained body is delivered exactly as the article's lines are written, so no omission receipt of any kind accompanies this package. Citations: the retained excerpts carry no citation command at all, so no bibliography entry of the article is transcribed and no reference is redistributed. Reference adaptation: none was needed and none was made; every reference inside the delivered unit resolves inside it, so no unresolved reference or citation is produced. Printed slips kept literal: no printed word, symbol, hypothesis or number of the counted statement or of its proof was altered. Layout and class: the article's own class is \texttt{amsart} with packages including inputenc, amsmath, amssymb, amsthm, booktabs, mathrsfs, enumerate, caption, accents, microtype, mathtools, longtable, varwidth, hyperref; the wrapper is the lane's pinned \texttt{amsart} editorial wrapper at 10pt on A4, which restates the theorem-like environments the retained text needs and adds the article's own macro definitions that the retained text needs (\texttt{\textbackslash EndsInf}, \texttt{\textbackslash HH}, \texttt{\textbackslash Hypers}, \texttt{\textbackslash MM}, \texttt{\textbackslash NN}), with extra wrapper packages (bm, bbm, dsfont, xspace, cancel, mathtools, enumitem). The delivered numbering is the unit's own. Assets: no retained span contains a figure environment, a graphics-inclusion command, a PDF-page inclusion, a table or a TikZ picture; no image file is copied into this package, the package declares no asset dependency and no image file is redistributed with it. Licence: version arXiv:2607.25065v1 of this article is distributed under the Creative Commons Attribution 4.0 International licence, as recorded in the arXiv licence metadata for this exact version and shown on the abstract page https://arxiv.org/abs/2607.25065v1. A full rights notice is delivered at licenses/arxiv-2607.25065-CC-BY-4.0.txt. Characters: every retained excerpt is pure ASCII, so the delivered engine, XeLaTeX with its default Latin Modern text font, reports no missing character.
\begin{lemma}[Mass transport principle]
\label{lemma:MTP}
A random point metric measure space $\mu$ is unimodular if and only if every Borel function $f : \MM^2 \to \left[0,\infty\right]$ satisfies
\begin{equation}
L_\mu(f) = R_\mu(f).
\end{equation}
\end{lemma}

\begin{cor}
\label{cor:thick part is infinitely ended}
Let $\mu$ be a unimodular random hyperbolic manifold with infinitely many infinite-volume ends. There is a sufficiently small  $\varepsilon_\mu > 0$ such that with positive  $\mu$-probability the pointed manifold $(M,p) \in \Hypers$ has $\mathrm{InjRad}_M(p) \ge \varepsilon_\mu$ and  the connected component of the $\varepsilon_\mu$-thick part containing the base point $p$ is an infinitely-ended manifold (with boundary).
\end{cor}

\begin{theorem}
\label{thm:random walk does not converge to a single end}
Let $\mu$ be a unimodular random hyperbolic manifold.  If $\mu$-almost every pointed manifold $(M,p)$ is transient and  satisfies $|\EndsInf{M}| > 1$ then for $\mu$-almost every pointed manifold $(M,p)$ the exit map $\zeta_{M,p} : T^1_p M \to \EndsInf{M}$  is not essentially constant.
\end{theorem}

\begin{proof}[Proof of Theorem \ref{thm:random walk does not converge to a single end} in the infinitely-ended case]
Let $\mu$ be unimodular random hyperbolic manifold. Assume towards contradiction that $\mu$-almost every pointed hyperbolic manifold $M$ is transient, satisfies $|\EndsInf{M}| = \infty$,  and admits a unique  infinite-volume end $\zeta_M \in \EndsInf{M}$ such that the geodesic flow trajectory starting at every base point and almost every unit tangent vector $(q,\vec u) \in T^1M$ accumulates to $\zeta_M$.

% Let $r > 0$ be a sufficiently large radius such that the ball $B_M(p,r)$  disconnects the random pointed manifold $(M,p)$ into at least two connected components with a positive $\mu$-probability. 
% Let $r>0$ be such that this probability is positive. 
For each $R > r > 0$ let $E_{r,R}$ be the following Borel set:
$$ E_{r,R} = \left\{ (M,x,y) \in \Hypers^{2} \: : \: 
\vcenter{
    \hbox{
      \begin{minipage}{7cm}
        \begin{enumerate}[label=\roman*)]
        \item $R \le d_M(x,y) \le R+1$,
          \item $M\setminus B(y,r)$ is not connected, and
          \item $x\in V$ for some unbounded connected component $V \in \mathcal{U}(B(y,r))$ which is \emph{not} a neighborhood of $\zeta_M$
        \end{enumerate}
      \end{minipage}
    }
  }\right\}.$$
The parameter $R$ will be chosen below (and will typically be much larger than $r$).

Assume that $(M,x_0,y_0) \in E_{r,R}$. We claim that any other point $y \in M$ with $(M,x_0,y) \in E_{r,R}$ must satisfy $d_M(y_0,y) < 2r+1$. Indeed, write $M \setminus B(y_0,r) = A \amalg B$ as a disjoint union of two collections of connected components, where $B$ is the connected end neighborhood of $\zeta_M$, and $A$ is the union of all other connected components. Namely  $A = M \setminus (B(y_0,r) \cup B)$ and $x_0 \in A$. There are several separate cases to consider:
\begin{itemize}
    \item Assume that $y \in B(y_0,r)$. Then the condition is clear, for certainly $d_M(y_0,y) \le r \le 2r+1$.
        \item Assume that $y \in A$. Let $V$ be the connected component of $M \setminus B_M(y,r)$ containing the point $x_0$. The condition $(M,x_0,y) \in E_{r,R}$ implies that $V$ is \emph{not} an end neighborhood of $\zeta_M$ and so the ball $B(y_0,r)$ cannot be entirely contained in $V$. Hence the ball $B(y,r)$ must intersect non-trivially either $B(y_0,r)$ or the geodesic arc $l$ from $x_0$ to $y_0$. In case $B(y,r) \cap B(y_0,r) \neq \emptyset$ the desired condition follows by the triangle inequality. In case $B(y,r) \cap l \neq \emptyset$ let $z$ be an arbitrary point of this intersection. Then $$d_M(x_0,y) \le d_M(x_0,z) + r.$$ 
    Since $d_M(x_0,y) \ge R$  we get $d_M(x_0,z) \ge R-r$. Since $d_M(x_0,y_0) \le R+1$ and as  $z$ lies on the geodesic arc from $x_0$ to $y_0$ we get $d_M(z,y_0) \le r+1$.  Therefore
    $$d_M(y_0,y) \le d_M(y_0,z) + d_M(z,y) \le (r+1)+r = 2r + 1$$
    as required.

    \item Assume that $y \in B$. The two points $x_0$ and $y$ lie in different connected components of $M \setminus B_M(y_0,r)$. Hence a geodesic path from $x_0$ to $y$ must intersect $B(y_0,r)$ at some point $y_1$.  We get
    $$ d_M(x_0,y_1) + d_M(y_1,y) = d_M(x_0,y) \le   R+ 1.$$ 
    The triangle inequality gives $d_M(x_0,y_1) \ge R-r$. Putting these two inequalities together gives $d_M(y,y_1) \le r+1$. Hence $d_M(y,y_0) \le 2r+1$ as required.

\end{itemize}

We now invoke the mass transport principle (Lemma \ref{lemma:MTP}). 
Let $F_{r,R} : \Hypers^2 \to \mathbb{R}$ be the characteristic function of the set $E_{r,R}$. 
The above claim gives the upper bound
$$ L_\mu (F_{r,R}) = \int_{\Hypers^2} F_{r,R}(M,x,y) \; \mathrm{dvol}_M(y) \; \mathrm{d}\mu(M,x) \le \mathrm{vol}_{\HH^n} B(2r+1).$$
We have used the fact that the hyperbolic volume of a ball in any hyperbolic manifold is upper bounded by the hyperbolic volume of the ball of the same radius in the hyperbolic space. The right-hand side is a constant depending only on $r$. 
We will arrive at a contradiction  by choosing $R$  to be sufficiently large, as follows. 

By Corollary \ref{cor:thick part is infinitely ended} there is an  $\varepsilon  >0$ such that the event $E$ defined by requiring that the pointed manifold $(M,p)$ has  $\mathrm{InjRad}_M(p) \ge \varepsilon$ and the connected component of the $\varepsilon$-thick part at the base point $p$ is infinitely-ended satisfies $\mu(E)>0$. 
As $r\to \infty$, the $\mu$-probability of the event $D$ that the ball $B_M(p,r)$ disconnects the random pointed manifold $(M,p)$ into at least two connected components tends to $1$.
Choose $r$ sufficiently large so that  $\mu(E\cap D) > 0$.

For each $N \in \NN$ we may find a radius $R_1(N) > 0$ such that the event $E_N$ given by requiring that $(M,p) \in E\cap D$ and that the space $M_{\ge \varepsilon} \setminus B(p,R_1(N))$ has at least $N$  connected components not accumulating to the end $\zeta_M$ satisfies $\mu(E_N) > \frac{1}{2}\mu(E\cap D)>0$.
Set $R = 
R(N) = R_1(N) + 2\varepsilon$. For each pointed manifold $(M,p) \in E_N$, it is possible to choose points $x_1,\ldots,x_N$ lying in each of those $N$ connected components. Note that $d_M(x_i,x_j) \ge 2\varepsilon$ for every  pair of distinct indices $i$ and $j$. Further $\mathrm{InjRad}_M(x_i) \ge \varepsilon$ for each $i$, since those points all lie in $M_{\ge \varepsilon}$. We obtain that $(M,z,p) \in E_{r,R}$ for all points $z \in \bigcup_{i=1}^N B(x_i,\varepsilon)$. This gives  the lower bound
$$ R_\mu(F_{r,R}) = \int_{\Hypers^2} F_{r,R}(M,x,y) \; \mathrm{dvol}_M(x) \; \mathrm{d}\mu(M,y)  \ge \mu(E_N)\cdot N \cdot \mathrm{vol}_{\HH^n} B(\varepsilon).$$
Note that $\mu(E_N) \cdot \mathrm{vol}_{\HH^n} B(\varepsilon)$ is bounded away from zero, while the parameter $N$ can be chosen arbitrarily large. We arrive at a contradiction to the mass transport principle by taking
 $$N \ge \frac{\mathrm{vol}_{\HH^n} B(2r+1)}{
 \tfrac12\mu(E\cap D) \cdot \mathrm{vol}_{\HH^n} B(\varepsilon)}$$
and choosing a sufficiently large radius $R = R(N)$ accordingly.
\end{proof}

\end{document}
